\section*{v4.6.0 Integrated Bounded-Readout and T2 Scope-Boundary Appendix}
\addcontentsline{toc}{section}{v4.6.0 Integrated Bounded-Readout and T2 Scope-Boundary Appendix}

\begin{definition}[v4.6.0 bounded readout surface]
The integrated v4.6.0 release-final bounded-readout surface is the finite guarded readout
surface
\[
\begin{aligned}
 W1&=\mathsf{pass},\\
 W2,W3,W4,W5_{\mathrm{formal\text{-}symbolic}},W6
   &=\mathsf{pass}_{\mathrm{bounded}},\\
 W7&=\mathsf{scope\text{-}boundary\ locked}.
\end{aligned}
\]
The status \(\mathsf{pass}_{\mathrm{bounded}}\) means that the corresponding
readout packet is admitted only in the declared bounded grammar.  It is not an
unrestricted all-readout witness.
\end{definition}

\begin{theorem}[Bounded readout closure in the declared grammar]
Inside the declared global-boundary grammar, the v4.6.0 bounded readout surface
is conservative and debt-free:
\[
 \mathsf{BRead}_{\mathrm{bounded}}
 \in
 \mathsf{Cons}(\mathcal G_{\mathrm{global}}^{\mathrm{decl}}),
 \qquad
 \Debt=0,
 \qquad
 \Leak=0.
\]
This theorem is a bounded formal-symbolic statement.  It does not promote an
unrestricted physical readout and does not assert an all-readout completion.
\end{theorem}

\begin{theorem}[T2 scope-boundary and non-implication]
The bounded readout theorem does not imply T2.  Equivalently,
\[
 \mathsf{BRead}_{\mathrm{bounded}}\not\Rightarrow T2,
 \qquad
 T2\in\mathsf{Ext}(\mathcal G_{\mathrm{global}}^{\mathrm{decl}}).
\]
The missing object is an all-readout witness family
\[
 \Theta^{\mathrm{all}}
 =
 \{\Theta_R:\mathcal G_{\mathrm{global}}^{\mathrm{decl}}\to R\}_{R\in\mathsf{Read}_{\mathrm{all}}}.
\]
No such family is constructed in this v4.6.0 release-final source.
\end{theorem}

\begin{definition}[Frozen all-readout obstruction]
The frozen v4.6.0 all-readout obstruction is
\[
\begin{aligned}
\mathfrak O_{\mathrm{all}}^{4.6.0}=\{&
\text{missing cofinal universal exterior selector},\\
&\text{missing all-readout naturality family},\\
&\text{unpromoted unrestricted exterior collar slot}\}.
\end{aligned}
\]
This obstruction is a scope boundary, not a contradiction.
\end{definition}

\begin{remark}[Millennium-matrix and Heisenberg-compensator status]
The v4.6.0 line separates the bounded, exterior, and mixed portions of the
Millennium-matrix and Heisenberg-compensator bookkeeping:
\[
\mathcal M=\mathcal M_{\mathrm{bd}}\oplus\mathcal M_{\mathrm{ext}}\oplus\mathcal M_{\mathrm{mix}},
\qquad
\mathcal H=\mathcal H_{\mathrm{bd}}\oplus\mathcal H_{\mathrm{ext}}\oplus\mathcal H_{\mathrm{mix}}.
\]
The bounded and finite exterior/mixed fragments are sufficient for the bounded
readout statement.  They are not sufficient for an all-readout theorem.
\end{remark}

\begin{definition}[Structure-first field-equation reset template]
The release-final field-equation template is
\[
\mathfrak E_{\mathrm{decl}}(\Phi;\mathcal M,\mathcal H,\mathsf{QAM})
=
\mathfrak G_{\mathrm{struct}}(\Phi)
+
\mathfrak C_{\mathcal H}(\Phi)
+
\mathfrak M_{\mathcal M}(\Phi)
+
\mathfrak Q_{\mathsf{QAM}}(\Phi)
-
\mathfrak J_{\mathrm{readout}}(\Phi).
\]
The bounded instance uses only the declared bounded and finite guarded fragments.
The all-readout instance is deliberately not asserted.
\end{definition}

\begin{remark}[Two-sided structural tensor notation]
The two-sided tensor notation is used only as a conservative compression layer.
Front or left-sided indices denote frame, readout, or language interfaces; back
or right-sided indices denote structure, action, or field interfaces.  Classical
upper--lower contraction remains ordinary action or readout.  Same-variance
comparisons require an explicitly declared comparison tensor and are not raw
Einstein contractions.
\end{remark}

\begin{remark}[v4.7 exterior--mixed notation guard]
The next MM/HC exterior--mixed field-equation program should write each proposed
mixed term with explicit sidedness.  Front-side indices carry readout, frame, or
language interfaces; back-side indices carry structure, action, or field
interfaces.  Any same-variance comparison between exterior and bounded channels
must pass through a declared comparison tensor or selector.  Thus the v4.7
program may sharpen \(\mathcal M_{\mathrm{ext}}\), \(\mathcal M_{\mathrm{mix}}\),
\(\mathcal H_{\mathrm{ext}}\), and \(\mathcal H_{\mathrm{mix}}\), but it may not
silently replace a comparison/readout distinction by an undeclared contraction.
\end{remark}

